% Standalone problem document, generated from the adjacent README.md.
% Compile with XeLaTeX. Mathematical content is maintained in Markdown.
\documentclass[11pt,a4paper]{article}
\usepackage[fontsize=10.5pt]{scrextend}
\usepackage[margin=22mm,headheight=15pt,headsep=9mm,footskip=11mm]{geometry}
\usepackage{fontspec}
\setmainfont{texgyrepagella}[Extension=.otf,UprightFont=*-regular,BoldFont=*-bold,ItalicFont=*-italic,BoldItalicFont=*-bolditalic]
\setsansfont{texgyreheros}[Extension=.otf,UprightFont=*-regular,BoldFont=*-bold,ItalicFont=*-italic,BoldItalicFont=*-bolditalic]
\setmonofont{texgyrecursor}[Extension=.otf,UprightFont=*-regular,BoldFont=*-bold,ItalicFont=*-italic,BoldItalicFont=*-bolditalic,Scale=MatchLowercase]
\usepackage{amsmath,amssymb}
\usepackage{unicode-math}
\setmathfont{texgyrepagella-math.otf}
\usepackage{xcolor,microtype,enumitem,needspace,fancyhdr}
\usepackage{bookmark}
\definecolor{ink}{HTML}{17344B}
\definecolor{accent}{HTML}{197D80}
\definecolor{muted}{HTML}{596773}
\hypersetup{colorlinks=true,linkcolor=accent,urlcolor=accent,pdftitle={MI-25 - Dimension-independent
Hlawka constants for Schatten
norms},pdfauthor={Open Problems in Numerical Linear Algebra}}
\urlstyle{same}
\setlength{\parindent}{0pt}
\setlength{\parskip}{5pt plus 1pt minus 1pt}
\linespread{1.04}
\setlength{\emergencystretch}{3em}
\widowpenalty=10000
\clubpenalty=10000
\displaywidowpenalty=10000
\allowdisplaybreaks[1]
\setlist{nosep,leftmargin=1.5em}
\providecommand{\tightlist}{\setlength{\itemsep}{0pt}\setlength{\parskip}{0pt}}
\providecommand{\pandocbounded}[1]{#1}
\pagestyle{fancy}
\fancyhf{}
\fancyhead[L]{\sffamily\footnotesize\color{muted}OPEN PROBLEMS IN NUMERICAL LINEAR ALGEBRA}
\fancyhead[R]{\sffamily\bfseries\footnotesize\color{accent}MI-25}
\fancyfoot[L]{\parbox[b]{0.9\textwidth}{\sffamily\fontsize{8}{10}\selectfont\color{muted}Editorial ratings; a bounded search cannot certify that a problem remains unsolved.\\Literature check: 2026-09-10}}
\fancyfoot[R]{\sffamily\footnotesize\color{muted}\thepage}
\renewcommand{\headrulewidth}{0.3pt}
\renewcommand{\headrule}{\hbox to\headwidth{\color{accent}\leaders\hrule height \headrulewidth\hfill}}
\makeatletter
\renewcommand{\section}{\Needspace{5\baselineskip}\@startsection{section}{1}{0pt}{12pt plus 2pt minus 2pt}{4pt}{\sffamily\bfseries\large\color{ink}}}
\renewcommand{\subsection}{\Needspace{4\baselineskip}\@startsection{subsection}{2}{0pt}{9pt plus 1pt minus 1pt}{3pt}{\sffamily\bfseries\normalsize\color{ink}}}
\makeatother
\setcounter{secnumdepth}{0}
\begin{document}
{\sffamily\small\color{muted}Matrix inequalities and norms\par}
\vspace{3pt}
{\raggedright\hyphenpenalty=10000\exhyphenpenalty=10000\sffamily\bfseries\fontsize{21}{24}\selectfont\color{ink}Dimension-independent
Hlawka constants for Schatten norms\par}
\vspace{7pt}
{\sffamily\small\textbf{Difficulty:} challenging\quad\textbf{Importance:} interesting
to the community\par}
{\sffamily\small\bfseries\color{accent}Status: Partially resolved\par}
\vspace{4pt}
{\color{accent}\hrule height 0.8pt}
\vspace{6pt}
\textbf{Rating rationale:} Sharp cancellation constants for entire
Schatten classes require new control across the Schatten scale, a
challenging norm-geometry question; dimension-independent control would
inform perturbation estimates across the community.

\subsection{Problem statement}\label{problem-statement}

For each \(1\le p<\infty\), determine the smallest \(C_p\in[0,\infty]\)
such that for every pair of positive integers \(N,M\) and every
\(X,Y,Z\in\mathbb C^{N\times M}\), \[
\|X\|_p+\|Y\|_p+\|Z\|_p-\|X+Y+Z\|_p
\le C_p\bigl(D_p(X,Y)+D_p(X,Z)+D_p(Y,Z)\bigr),
\] where \(D_p(U,V)=\|U\|_p+\|V\|_p-\|U+V\|_p\) and
\(\|U\|_p=(\sum_j s_j(U)^p)^{1/p}\). The value \(C_p=\infty\) means that
no finite constant works uniformly over dimensions. The inequality must
also hold when the parenthesized expression vanishes; a ratio definition
must not discard such triples.

This asks whether the cancellation loss in a three-term matrix sum can
be controlled sharply by the cancellation losses of its pairs. Such
estimates concern the geometry of the matrix norms used in perturbation
bounds and low-rank approximation.

\subsection{References}\label{references}

\begin{enumerate}
\def\labelenumi{\arabic{enumi}.}
\tightlist
\item
  K. M. R. Audenaert and F. Kittaneh, \emph{Problems and Conjectures in
  Matrix and Operator Inequalities}, arXiv:1201.5232v3 (2012), §8.2,
  Problem 7, equation (48), and the following discussion.
  \href{https://arxiv.org/html/1201.5232}{Full text}.
\item
  Audenaert and Kittaneh, published version in \emph{Études
  opératorielles}, Banach Center Publications 112 (2017), §8.2, Problem
  6, equation (48), p.~25.
  \href{https://www.impan.pl/shop/en/publication/transaction/download/product/111309}{Publisher
  PDF}.
\end{enumerate}

Status check (2026-09-10): the source gives \(C_2=1\) and demonstrates
nonexistence of a finite constant for the operator norm, while posing
the Schatten constants generally. Current source records and searches
for ``Hlawka Schatten constant'', ``Hlawka operators Audenaert
Kittaneh'', and 2025/2026 found no determination of all \(C_p\). Later
determinantal Hlawka inequalities concern a different inequality and do
not answer this norm-deficit question. The search is bounded; individual
\(p\)-cases should be rechecked before treating them as new results.
\end{document}
